\slajdid{08-s008}
\begin{frame}[label=08-s008]{Decoupling capacitors}
\fontsize{9}{10}\selectfont\setlength{\parskip}{1.5pt}
Kondenzatori za blokadu\par
\begin{columns}[c,onlytextwidth]\column{.65\textwidth}\input{slike/tikz/s008_dekapling_potrosac.tex}
\column{.33\textwidth}
$V_{CC1}(t_0^-)=\dfrac{1000}{1001}V_{CC}$\par\medskip
$V_{CC1}(t_0^+)=\dfrac{1000}{1001}V_{CC}$
\end{columns}
\begin{columns}[T,onlytextwidth]\column{.55\textwidth}
Od „svoje potrošnje“ - Transient load $C_{TL}$\par
Od šuma koji je posledica rada drugih kola i smetnji - Noise $C_N$\par
Isti efekat – jedan kondenzator uz svako kolo\par
\centering\begin{tikzpicture}[x=.85cm,y=.65cm,font=\sitneoznake,>=Latex]
\draw (-0.22,0.22)--(0.24,0)--(-0.22,-0.22)--cycle;
\draw (0.29,0)circle(.05);\node at(-0.08,0){1};
\draw(-0.7,0)--(-0.22,0);\draw(0.34,0)--(0.85,0);
\draw(-0.1,.16)--(-0.1,.8);\draw(-0.1,-.16)--(-0.1,-1.05);
\draw(-0.1,.6)--(-1.3,.6)--(-1.3,.05);\draw(-1.3,-.05)--(-1.3,-.6)--(-0.1,-.6);
\draw(-1.5,.05)--(-1.1,.05);\draw(-1.5,-.05)--(-1.1,-.05);\node[left]at(-1.5,0){$C$};
\draw[->,DEcrvena](-0.25,-.4)--(-0.25,-.48)--(-1.03,-.48)--(-1.03,.48)--(-0.25,.48)--(-0.25,.35);
\draw (3.1799999999999997,0.22)--(3.6399999999999997,0)--(3.1799999999999997,-0.22)--cycle;
\draw (3.69,0)circle(.05);\node at(3.32,0){2};
\draw(2.7,0)--(3.1799999999999997,0);\draw(3.7399999999999998,0)--(4.25,0);
\draw(3.3,.16)--(3.3,.8);\draw(3.3,-.16)--(3.3,-1.05);
\draw(3.3,.6)--(2.0999999999999996,.6)--(2.0999999999999996,.05);\draw(2.0999999999999996,-.05)--(2.0999999999999996,-.6)--(3.3,-.6);
\draw(1.9,.05)--(2.3,.05);\draw(1.9,-.05)--(2.3,-.05);\node[left]at(1.9,0){$C$};
\draw[->,DEcrvena](3.15,-.4)--(3.15,-.48)--(2.37,-.48)--(2.37,.48)--(3.15,.48)--(3.15,.35);
\draw(-2.8,.8)node[above]{$V_{CC}$}--(-2.4,.8);\draw(-2.4,.7)rectangle(-1.8,.9);\draw(-1.8,.8)--(1.4,.8);\draw(1.4,.72)rectangle(1.75,.88);\draw(1.75,.8)--(4.15,.8);\draw(4.15,.68)rectangle(5.2,.92);\draw(5.2,.8)--(5.5,.8);
\draw(-2.8,-1.05)node[ground]{}--(-1.9,-1.05);\draw(-1.9,-1.15)rectangle(-1.75,-.95);\draw(-1.75,-1.05)--(1.2,-1.05);\draw(1.2,-1.15)rectangle(1.8,-.95);\draw(1.8,-1.05)--(4.75,-1.05);\draw(4.75,-1.12)rectangle(5.0,-.98);\draw(5,-1.05)--(5.5,-1.05);
\draw[->](-2.1,1)--(-.9,1)node[midway,above]{$i_{cc}(t)$};
\draw[->](.4,-1.3)--(-.8,-1.3)node[midway,below]{$i_{DGND}(t)$};
\node at(5,0){$\cdots$};
\end{tikzpicture}
\par
{\sitneoznake «različite kockice, različite impedanse i parazitni efekti»\;
\tikz[baseline=-.5ex]{\draw(0,0)circle(.09);\fill(-.03,.025)circle(.008);\fill(.03,.025)circle(.008);\draw(-.045,-.04)..controls(0,0)..(.045,-.04);}}
\column{.43\textwidth}
Kada se ponovo potrošač vrati na 1000R, kondenzator će se dopuniti iz izvora za napajanje. „Privremena baterija“\par
Najčešći dekapling kondenzator uz logička kola je keramički, vrednosti 100nF.\par
Većina digitalnih sistema ne bi ispravno radila da ne postoje ovi kondenzatori. Česta greška mladih inženjera jeste da ih ne postave, ili ne postave na odgovarajuća mesta, uz samo napajanje čipova.\par
Linije između logičkog kola i dekapling kondenzatora treba da su što kraće zbog parazitnih efekata
\end{columns}
\end{frame}
